%% VisualReF — Master's Thesis Draft
%% Author: R. Boghean, Maastricht University
%% Style: ACM-inspired academic research paper (single-column, 12pt, double-spaced)
%%
%% Compile with:
%%   pdflatex thesis.tex && bibtex thesis && pdflatex thesis.tex && pdflatex thesis.tex

\documentclass[12pt,a4paper]{article}

%% ── Packages ──────────────────────────────────────────────────────────────────
\usepackage[utf8]{inputenc}
\usepackage[T1]{fontenc}
\usepackage[british]{babel}
\usepackage{lmodern}
\usepackage{microtype}

\usepackage[
  top=2.5cm, bottom=2.5cm,
  left=3cm,  right=3cm,
]{geometry}

\usepackage{setspace}
\onehalfspacing

\usepackage{amsmath,amssymb,amsfonts}
\usepackage{bm}
\usepackage{booktabs}
\usepackage{multirow}
\usepackage{array}
\usepackage{tabularx}
\usepackage{graphicx}
\usepackage{xcolor}
\usepackage{url}
\usepackage{hyperref}
\usepackage{cleveref}
\usepackage{algorithm}
\usepackage{algpseudocode}
\usepackage{caption}
\usepackage{subcaption}
\usepackage{enumitem}
\usepackage{parskip}
\usepackage{titlesec}
\usepackage{fancyhdr}
\usepackage{natbib}

%% ── Hyperref config ───────────────────────────────────────────────────────────
\hypersetup{
  colorlinks = true,
  linkcolor  = blue!60!black,
  citecolor  = blue!60!black,
  urlcolor   = blue!60!black,
  pdftitle   = {VisualReF: Multimodal Relevance Feedback for Interactive Image Retrieval},
  pdfauthor  = {R. Boghean},
}

%% ── Header / footer ──────────────────────────────────────────────────────────
\pagestyle{fancy}
\fancyhf{}
\rhead{\small VisualReF: Multimodal Relevance Feedback}
\lhead{\small R.\ Boghean --- Maastricht University}
\cfoot{\thepage}
\renewcommand{\headrulewidth}{0.4pt}

%% ── Section styling ──────────────────────────────────────────────────────────
\titleformat{\section}{\large\bfseries}{\thesection.}{0.5em}{}
\titleformat{\subsection}{\normalsize\bfseries}{\thesubsection.}{0.5em}{}
\titleformat{\subsubsection}{\normalsize\itshape}{\thesubsubsection.}{0.5em}{}

%% ── Math macros ──────────────────────────────────────────────────────────────
\newcommand{\vq}{\bm{q}}
\newcommand{\vd}{\bm{d}}
\newcommand{\ve}{\bm{e}}
\newcommand{\norm}[1]{\left\lVert#1\right\rVert}

%% ── Metadata ──────────────────────────────────────────────────────────────────
\title{%
  \textbf{VisualReF: Multimodal Relevance Feedback\\
  for Interactive Image Retrieval}\\[0.4em]
  \large A Prototype Integrating SAM~3 Segmentation,
  SigLIP Embeddings, and Vision-Language Models
}

\author{%
  Rares Boghean\\[0.2em]
  \small Department of Advanced Computing Sciences\\
  \small Maastricht University, The Netherlands\\
  \small \texttt{r.boghean@student.maastrichtuniversity.nl}
}

\date{April 2026 --- Master's Thesis Draft}

%% =============================================================================
\begin{document}

\maketitle
\thispagestyle{fancy}

%% ── Abstract ──────────────────────────────────────────────────────────────────
\begin{abstract}
\noindent
Interactive image retrieval systems must bridge the semantic gap between
user intent—expressed visually through pointing and selection—and the
numerical representations that underpin modern nearest-neighbour search.
We present \textbf{VisualReF}, an end-to-end prototype that closes this gap
by combining three complementary technologies: (1)~the SigLIP vision-language
encoder for a shared image-text embedding space, (2)~the SAM~3 interactive
segmenter for pixel-precise region selection, and (3)~the Llama~3.2-Vision
large language model for automatic, query-aware captioning of selected regions.
User feedback—both visual (click-to-segment) and textual (free-form hints)—
is fused into a Rocchio-style query update that incrementally shifts the search
query towards relevant and away from irrelevant content.
On the Visual Genome corpus (\(\sim\)108\,000 images) with a FAISS flat
inner-product index, the prototype demonstrates interactive retrieval at a
resolution of 1024-dimensional SigLIP embeddings.
We describe the system architecture in full, formalise the multimodal Rocchio
update, identify the key engineering challenges—in particular the latency
introduced by large vision-language model inference—and outline a structured
programme of experiments designed to evaluate retrieval effectiveness,
multimodal fusion strategies, and scalability.
The prototype is intended as the empirical basis for a forthcoming evaluation
study submitted to the ACM Recommender Systems conference (RecSys~'25).
\end{abstract}

\noindent\textbf{Keywords:} interactive image retrieval, relevance feedback,
vision-language models, Rocchio algorithm, segment anything model,
multimodal fusion, FAISS, SigLIP.

\tableofcontents
\clearpage

%% =============================================================================
\section{Introduction}
\label{sec:introduction}

Content-based image retrieval (CBIR) has undergone a paradigm shift in recent
years. Early systems represented images using hand-crafted descriptors—colour
histograms, SIFT keypoints, local binary patterns—and measured similarity in
that hand-crafted space~\citep{Smeulders2000}. Contemporary systems instead
project images and text jointly into a shared high-dimensional embedding space
learned from hundreds of millions of image-caption pairs, with models such as
CLIP~\citep{Radford2021} and SigLIP~\citep{Zhai2023} providing the state of
the art. Given such a shared space, a natural language query can be embedded
alongside the image collection, enabling semantic retrieval with no manual
feature engineering.

Yet even with learned embeddings, the \emph{semantic gap}~\citep{Smeulders2000}
persists: the gap between what a user \emph{intends}—a precise visual concept
that may be difficult to articulate in words—and what a text query can
communicate. A user searching for \emph{``a dog playing fetch near water at
sunset''} may not articulate every constraint in advance; they recognise
relevant images when they see them. Relevance feedback~\citep{Rocchio1971}
addresses this by treating retrieval as an iterative dialogue: the user
inspects an initial result set, labels some images as relevant or irrelevant,
and the system updates the query accordingly.

Classical relevance feedback operates on \emph{whole images}: the user selects
entire images as relevant or irrelevant, and the system moves the query vector
towards relevant images and away from irrelevant ones in embedding space.
This coarse granularity is a limitation. A returned image may contain both
the desired object (a dog) and distracting context (a car park). Labelling
the whole image as relevant inadvertently reinforces the unwanted content.

Interactive segmentation offers a solution. The Segment Anything Model
(SAM)~\citep{Kirillov2023} and its successors SAM~2~\citep{Ravi2024} and
SAM~3~\citep{Meta2025} allow users to click a single point on an object and
obtain a pixel-accurate mask for that object in milliseconds.
If relevance feedback operates on \emph{segmented regions} rather than
whole images, the feedback signal can be made region-specific: only the
pixels the user actually cares about contribute to the query update.

A further limitation of image-only feedback is representational: SigLIP image
embeddings capture the global visual content of a crop but may miss the
semantic concept the user has in mind. A vision-language model (VLM) such as
Llama~3.2-Vision~\citep{Meta2024} can describe the content of a selected
region in natural language, and the resulting caption can be embedded as text
in the same SigLIP space. Since SigLIP is trained to align image and text
representations, a textual description of the selected region provides a
complementary and often more discriminative signal than the image embedding
alone.

\textbf{VisualReF} integrates these three components—SigLIP, SAM~3, and a
VLM—into a single interactive retrieval prototype. The system contributions
are as follows:

\begin{enumerate}[leftmargin=1.5em]
  \item A multimodal Rocchio update that fuses (i)~image embeddings of SAM
        crops, (ii)~textual embeddings of automatically generated VLM captions,
        and (iii)~free-form user text hints, with a weighted combination whose
        proportions adapt to signal availability.
  \item An asynchronous captioning pipeline in which VLM inference is triggered
        immediately after segmentation and cached, so that captions are
        available with negligible additional latency at feedback time.
  \item Integration with the Visual Genome region description
        corpus~\citep{Krishna2017}, providing instant, human-written region
        phrases as a fallback when VLM inference is unavailable.
  \item A full-stack open prototype (FastAPI back-end, Next.js front-end)
        suitable for controlled user studies.
\end{enumerate}

The remainder of this thesis is organised as follows.
\Cref{sec:related} reviews related work on content-based image retrieval,
relevance feedback, visual segmentation, and vision-language models.
\Cref{sec:system} describes the system architecture.
\Cref{sec:method} formalises the multimodal Rocchio update and the signal
fusion strategy.
\Cref{sec:impl} details key implementation decisions.
\Cref{sec:experiments} proposes an experimental programme and evaluation
methodology.
\Cref{sec:discussion} discusses limitations, open questions, and future work.
\Cref{sec:conclusion} concludes.

%% =============================================================================
\section{Related Work}
\label{sec:related}

\subsection{Content-Based Image Retrieval}

CBIR research dates to the early 1990s with systems such as QBIC~\citep{Flickner1995}
and VisualSEEk~\citep{Smith1996}. These systems retrieved images based on
low-level visual features—colour, texture, shape—and suffered acutely from the
semantic gap: visual similarity in feature space does not correspond to
semantic similarity as perceived by users.

The introduction of deep convolutional neural networks brought a step change
in retrieval quality~\citep{Babenko2014,Tolias2016}. Features extracted from
intermediate CNN layers served as compact image descriptors, and approximate
nearest-neighbour structures such as
FAISS~\citep{Johnson2019} made billion-scale retrieval practical.

The most recent generation of retrieval systems is built on contrastive
vision-language models. CLIP~\citep{Radford2021} and SigLIP~\citep{Zhai2023}
learn a joint embedding space in which semantically related images and text
captions cluster together. This enables \emph{zero-shot} retrieval: an
arbitrary natural-language query can be embedded and used to retrieve relevant
images with no per-category training.

\subsection{Relevance Feedback}

The Rocchio algorithm~\citep{Rocchio1971} was originally developed for text
retrieval in the vector space model. Given a query vector $\vq$ and sets of
relevant and irrelevant documents $D^+$ and $D^-$, the updated query is:
\begin{equation}
  \vq' = \alpha\,\vq
       + \beta\,\frac{1}{|D^+|}\sum_{\vd \in D^+} \vd
       - \gamma\,\frac{1}{|D^-|}\sum_{\vd \in D^-} \vd.
  \label{eq:rocchio_text}
\end{equation}
\citet{Rui1998} extended this formulation to image retrieval, demonstrating
that relevance feedback over visual descriptors substantially improves
retrieval precision in iterative sessions.

Subsequent work explored support vector machine (SVM)-based
feedback~\citep{Tong2001}, active learning
strategies~\citep{Tong2001active}, and kernel-based
methods~\citep{Gosselin2004}. A comprehensive survey appears
in~\citet{Zhou2003}.

In the deep learning era, relevance feedback has been revisited with
neural query reformulation models~\citep{Jiang2015} and attention-based
weighting of relevant features~\citep{Yu2021}. VisualReF revisits the
classical Rocchio update but applies it in the unified embedding space of a
modern vision-language model, where image and text signals can be combined
without modality-specific adaptation.

\subsection{Interactive Segmentation}

Interactive segmentation traces to approaches based on graph cuts~\citep{Boykov2001}
and GrabCut~\citep{Rother2004}, which required significant user annotation
effort. Deep interactive segmentation methods~\citep{Xu2016,Sofiiuk2022}
reduced annotation burden substantially.

SAM~\citep{Kirillov2023} is a promptable segmentation model trained on a
dataset of over one billion masks. Given a single point click, a bounding box,
or a free-form text prompt, SAM produces high-quality binary masks across a
wide range of object categories without per-class fine-tuning. SAM~2~\citep{Ravi2024}
extended the model to video through a streaming memory mechanism. SAM~3~\citep{Meta2025}
incorporates a vision-language backbone that enables text-guided segmentation
and interactive instance refinement, which VisualReF exploits for per-click
mask prediction.

\subsection{Vision-Language Models for Retrieval and Description}

The alignment between vision and language achieved by models such as
CLIP~\citep{Radford2021} and ALIGN~\citep{Jia2021} established the foundation
for multimodal retrieval. LLaVA~\citep{Liu2023} demonstrated that large
language models augmented with a vision encoder can produce detailed, accurate
descriptions of image content following natural-language instructions.

Several systems have investigated using VLM-generated descriptions to augment
retrieval. \citet{Luo2023} showed that LLM-generated query expansions improve
text-to-image retrieval; VisualReF applies an analogous idea to \emph{query
feedback}: rather than expanding the initial text query, the VLM expands the
\emph{feedback signal} by describing the visual content the user indicated as
relevant or irrelevant.

\subsection{Visual Genome and Region Descriptions}

The Visual Genome dataset~\citep{Krishna2017} contains 108,249 images
annotated with region descriptions—short natural-language phrases (``a man
wearing a blue shirt'') linked to bounding boxes. VisualReF leverages these
descriptions as a source of instant, human-authored region phrases, indexed
by image identifier and retrieved by intersection-over-union (IoU) matching
between the SAM mask and annotated bounding boxes.

\subsection{Positioning VisualReF}

VisualReF is most closely related to interactive retrieval systems that combine
visual feedback with semantic description. \citet{Custodio2023} explored using
captioned image regions for query refinement; the prior VisualReF
demonstration paper by \citet{Khaertdinov2025} introduced the combined SAM
+ Rocchio + SigLIP pipeline. The present thesis extends that baseline with
(i)~the integration of VLM-generated captions as a third feedback signal,
(ii)~an asynchronous captioning architecture that hides VLM latency, and
(iii)~a systematic experimental programme to evaluate the contribution of
each signal.

%% =============================================================================
\section{System Architecture}
\label{sec:system}

\begin{figure}[t]
  \centering
  % Replace with actual figure when available
  \fbox{\parbox{0.92\textwidth}{\centering\vspace{2cm}
    \textit{[Architecture diagram: browser UI $\leftrightarrow$ FastAPI back-end.
    Three back-end subsystems: SigLIP/FAISS, SAM~3, Ollama/Llama~3.2-Vision.
    VG region index shown as fourth data source.]}
  \vspace{2cm}}}
  \caption{High-level architecture of VisualReF.
  The browser communicates directly with the FastAPI back-end (port~8001).
  The back-end integrates four subsystems: a SigLIP encoder with a FAISS index
  for retrieval, SAM~3 for region segmentation, Ollama for VLM inference, and
  the Visual Genome region description index for instant textual phrases.}
  \label{fig:architecture}
\end{figure}

The VisualReF prototype follows a two-tier client-server architecture
(\cref{fig:architecture}). The front-end is a Next.js 15 single-page
application; the back-end is a FastAPI server that hosts all model inference.
To avoid proxy timeout issues inherent in the Next.js development server,
the browser calls the FastAPI server directly via a configurable base URL.

\subsection{Back-End Components}

\paragraph{SigLIP Encoder.}
We use \texttt{google/siglip-large-patch16-256}~\citep{Zhai2023}, a 400M
parameter vision-language model that maps both images and text into a shared
1024-dimensional $\ell_2$-normalised embedding space. Unlike CLIP, which uses
a softmax contrastive loss, SigLIP employs a sigmoid binary cross-entropy loss
over image-text pairs, which has been shown to improve alignment quality,
particularly for fine-grained text queries~\citep{Zhai2023}.

Images are processed at $256 \times 256$ pixels with 16-pixel patch size
(hence the 256 patch tokens per image). Text queries are tokenised with the
SigLIP vocabulary (max 64 tokens). All inference runs under
\texttt{torch.no\_grad()} with the model in evaluation mode.

\paragraph{FAISS Index.}
The image collection is indexed offline: each image is resized to $256\times
256$ pixels, encoded with the SigLIP image encoder, $\ell_2$-normalised, and
stored in a \texttt{faiss.IndexFlatIP} index (exact cosine similarity via inner
product on unit vectors). The index for Visual Genome contains 108,249 vectors
of dimension 1024, occupying approximately 420~MB on disk.

At query time, the text embedding of the user's query is used directly as the
FAISS search vector. Because the vectors are unit-normalised, inner product
is equivalent to cosine similarity.

\paragraph{SAM~3 Segmenter.}
SAM~3~\citep{Meta2025} is loaded in interactive image mode using
\texttt{facebook/sam3} from the HuggingFace Hub with
\texttt{enable\_inst\_interactivity=True}, which activates the dual-backbone
architecture that produces both SAM-3-native and SAM-2-compatible feature
levels. Concretely, this triggers the \texttt{Sam3DualViTDetNeck} with
\texttt{add\_sam2\_neck=True}, producing three feature levels at spatial
resolutions $\{288^2, 144^2, 72^2\}$, which are consumed by the SAM~2
memory-based interactive predictor.

For each user click event, the system:
\begin{enumerate}[leftmargin=1.5em,noitemsep]
  \item Scales click coordinates from browser CSS pixels to the full-resolution
        image pixel space.
  \item Computes a bounding box prompt from the click ensemble
        (padded at 20\% for foreground-only prompts, 13\% for negative-only).
  \item Calls \texttt{SAM3InteractiveImagePredictor.predict()} with
        multimask output for single-click prompts and single-mask output
        for multi-click prompts.
  \item Selects the best mask by score, subject to a maximum area fraction
        of 28\% of the image (to reject trivially large masks).
  \item Downsamples the full-resolution boolean mask to the preview resolution
        (equal to the browser's rendered image dimensions) and encodes it as
        run-length encoding (RLE) in COCO convention.
\end{enumerate}
All SAM inference runs under \texttt{torch.no\_grad()} to prevent autograd
graph accumulation, which is critical for correctness on Apple MPS where
several autograd operations are unsupported.

\paragraph{VLM Captioning.}
Captions are generated by Llama~3.2-Vision (11B parameters) served via
Ollama~\citep{Meta2024}. Given a SAM-cropped region image and the current
search query, the model is prompted to produce a concise (15-word) description
that focuses on the aspect of the crop most relevant to the query.
The prompt provides both the isolated crop and a full-scene context image
(with a bounding box highlight) so the VLM can ground the description in
spatial context. Inference temperature is set to 0.1 (low creativity) and
maximum output tokens to 60.

A key engineering contribution is the \emph{asynchronous caption pipeline}:
immediately after the \texttt{/segment} endpoint returns the RLE mask, the
server fires a background task (\texttt{asyncio.create\_task}) that calls
the VLM on the segmented crop. The resulting caption is stored in a bounded
LRU cache keyed by \texttt{(image\_path, label, query, user\_hint)}.
When the user subsequently submits feedback via \texttt{/apply\_feedback},
the caption is retrieved from the cache with zero additional latency.

\paragraph{Visual Genome Region Index.}
For the Visual Genome corpus, a runtime phrase index is constructed from the
\texttt{region\_descriptions.json} file. For each retrieved image, the
index returns region phrases whose ground-truth bounding boxes have the highest
IoU with the SAM mask. These phrases are returned immediately with the
\texttt{/segment} response and reused during feedback, bypassing VLM inference
entirely for the majority of images in the corpus.

\subsection{Front-End}

The front-end is built with Next.js~15 and React~19, using Zustand for global
state management. The principal UI components are:

\begin{itemize}[leftmargin=1.5em,noitemsep]
  \item \textbf{ImageGallery / ImageCard}: renders the top-$k$ result images.
        Each card hosts an HTML5 Canvas overlay that renders SAM masks in
        real-time. Click coordinates are captured relative to the card's
        rendered bounding rectangle and converted to the image's natural pixel
        space before being sent to the server.
  \item \textbf{FeedbackPanel}: collects user text hints (``more of / less of'')
        and displays AI Vision descriptions auto-filled from VLM captions,
        with a one-click ``Use'' button to transfer a caption to the hint field.
\end{itemize}

The mask overlay is rendered as follows: the RLE-encoded mask is decoded on the
client to a flat \texttt{Uint8ClampedArray}, painted onto an off-screen canvas,
then composited at 40\% opacity onto the image display rectangle (not the full
canvas) using \texttt{drawImage} with source and destination coordinates,
ensuring correct spatial alignment regardless of image scaling.

\subsection{API Endpoints}

\begin{table}[h]
\centering
\caption{HTTP API exposed by the FastAPI back-end.}
\label{tab:api}
\begin{tabularx}{\textwidth}{@{}llX@{}}
\toprule
\textbf{Method} & \textbf{Path} & \textbf{Description} \\
\midrule
POST & \texttt{/search}          & Text query $\to$ top-$k$ images (base64 + paths + scores) \\
POST & \texttt{/segment}         & Point prompts $\to$ SAM mask (RLE), region crop, VG phrases, cached caption \\
POST & \texttt{/apply\_feedback} & Annotated images $\to$ Rocchio update $\to$ new top-$k$ images \\
POST & \texttt{/caption}         & On-demand Ollama captioning of a base64 region \\
GET  & \texttt{/health}          & Server liveness; model and index status \\
GET  & \texttt{/sam\_status}     & SAM~3 load status \\
GET  & \texttt{/ollama\_status}  & Ollama availability and active model \\
GET  & \texttt{/metrics}         & CPU, memory, GPU, FAISS index size (TTL-cached) \\
\bottomrule
\end{tabularx}
\end{table}

%% =============================================================================
\section{Methodology: Multimodal Rocchio Feedback}
\label{sec:method}

\subsection{Problem Formulation}

Let $\mathcal{I} = \{I_1, \ldots, I_N\}$ be a corpus of $N$ images, each
encoded into a unit-normalised embedding $\ve_i \in \mathbb{R}^d$ ($d=1024$).
The retrieval task is to find the $k$ images most similar to the user's
information need, expressed as a query vector $\vq \in \mathbb{R}^d$.

At round $t$, the user is shown the top-$k$ results for $\vq^{(t)}$. They
interact with a subset of result images by:
\begin{enumerate}[leftmargin=1.5em,noitemsep]
  \item Clicking on one or more objects, producing SAM masks.
  \item Optionally typing free-text hints.
\end{enumerate}
From these interactions the system constructs a \emph{positive signal}
$\vd^+$ and a \emph{negative signal} $\vd^-$, both in $\mathbb{R}^d$, and
computes an updated query $\vq^{(t+1)}$.

\subsection{Signal Extraction}

\paragraph{Image signal from SAM crops.}
For each annotated mask $m$ on image $I$, the mask is upsampled from preview
resolution to the original image resolution via nearest-neighbour interpolation.
The region $r = \text{crop}(I, m)$ is extracted by setting non-masked pixels
to neutral grey (RGB value 128), then cropping to the tight bounding box of
the mask. This ensures SigLIP patch embeddings are not biased by background
content outside the selected region.

The crop is encoded by the SigLIP image encoder:
\begin{equation}
  \ve^{\text{img}}_m = f_{\text{img}}(r) \in \mathbb{R}^d.
\end{equation}
When there are multiple relevant (resp.\ irrelevant) crops, the positive
(resp.\ negative) image signal is their mean embedding:
\begin{equation}
  \vd^+_{\text{img}} = \frac{1}{|M^+|}\sum_{m \in M^+} \ve^{\text{img}}_m,
  \qquad
  \vd^-_{\text{img}} = \frac{1}{|M^-|}\sum_{m \in M^-} \ve^{\text{img}}_m.
\end{equation}

\paragraph{Text signal from VLM captions, VG phrases, and user hints.}
Three sources of text contribute to the text signal:
\begin{itemize}[leftmargin=1.5em,noitemsep]
  \item $C_{\text{VLM}}$: Llama~3.2-Vision captions of SAM crops (from cache).
  \item $P_{\text{VG}}$: Visual Genome region phrases matched by IoU.
  \item $H$: Free-form user hints (one positive string, one negative string).
\end{itemize}
These are concatenated into a positive text list $T^+$ and a negative text
list $T^-$. Duplicates are removed and the list is capped at 20 entries to
bound inference time. The positive text signal is:
\begin{equation}
  \vd^+_{\text{txt}} = \frac{1}{|T^+|}\sum_{s \in T^+} f_{\text{txt}}(s),
\end{equation}
where $f_{\text{txt}}(s)$ is the SigLIP text encoder applied to string $s$.

\paragraph{VG phrase priority.}
When VG phrases are available and no user hint is provided, VLM captioning
is skipped for that polarity—VG phrases already provide rich, grounded
descriptions at zero latency.

\subsection{Multimodal Fusion}

The image and text signals are combined by weighted addition:
\begin{equation}
  \vd^+ = w_{\text{img}}\,\vd^+_{\text{img}} + w_{\text{txt}}\,\vd^+_{\text{txt}},
  \label{eq:fusion}
\end{equation}
and analogously for $\vd^-$. The weights adapt to signal availability:
\begin{equation}
  (w_{\text{img}},\, w_{\text{txt}}) =
  \begin{cases}
    (0.4,\, 0.6) & \text{if VLM captions are available,} \\
    (0.5,\, 0.5) & \text{otherwise.}
  \end{cases}
\end{equation}
The rationale is that text embeddings in the SigLIP space carry higher
semantic specificity than image embeddings when the text describes the
relevant concept explicitly; the higher text weight when VLM captions are
available reflects this.
If one signal is absent (no crops, or no text), only the present signal is used.

\subsection{Rocchio Update}

The query is updated according to the Rocchio formula adapted to the
multimodal setting:
\begin{equation}
  \vq^{(t+1)} = \ell_2\text{-normalise}\!\left(
    \alpha\,\vq^{(t)} + \beta\,\vd^+ - \gamma\,\vd^-
  \right),
  \label{eq:rocchio}
\end{equation}
with parameters $\alpha = 0.8$, $\beta = 0.5$, $\gamma = 0.15$.
The high $\alpha$ preserves the original search intent; the asymmetry
$\beta > \gamma$ reflects the practical observation that users tend to provide
more positive than negative feedback and that over-penalising irrelevant
content risks over-correction.

\paragraph{Query drift mitigation.}
Across multiple feedback rounds, the accumulated query may drift away from the
user's original intent as feedback signals compound. An optional
\emph{fuse-initial-query} mechanism anchors the update:
\begin{equation}
  \vq^{(t)}_{\text{fused}} = \frac{\vq^{(t)} + \vq^{(0)}}{2},
\end{equation}
where $\vq^{(0)}$ is the embedding of the initial text query. The Rocchio
update is applied to $\vq^{(t)}_{\text{fused}}$ rather than $\vq^{(t)}$.
This is user-controllable via a checkbox in the UI.

\paragraph{Null signal handling.}
When a feedback round produces neither positive nor negative signals
(e.g.\ the user provides a text query only), the Rocchio update degenerates
to $\vq^{(t+1)} = \ell_2\text{-normalise}(\alpha\,\vq^{(t)})$, equivalent
to a new search with the most recently accumulated query.

\subsection{Retrieval}

The updated query $\vq^{(t+1)}$ is passed to FAISS:
\begin{equation}
  \text{Results} = \operatorname*{arg\,top\text{-}k}_{i \in \{1,\ldots,N\}}
    \langle \vq^{(t+1)},\, \ve_i \rangle.
\end{equation}
Because both query and index vectors are unit-normalised, this is equivalent
to cosine similarity search.

\begin{algorithm}[t]
\caption{VisualReF feedback round}
\label{alg:feedback}
\begin{algorithmic}[1]
\Require Current query $\vq^{(t)}$, user annotations $\mathcal{A}$,
         VLM caption cache $C$, VG phrase index $P$
\Ensure Updated results $\mathcal{R}^{(t+1)}$
\State $M^+ \gets \{m \in \mathcal{A} : \text{label}(m) = \text{Relevant}\}$
\State $M^- \gets \{m \in \mathcal{A} : \text{label}(m) = \text{Irrelevant}\}$
\State $\vd^+_{\text{img}} \gets \text{MeanEmbed}(\{\text{crop}(I, m) : m \in M^+\})$
\State $\vd^-_{\text{img}} \gets \text{MeanEmbed}(\{\text{crop}(I, m) : m \in M^-\})$
\State $T^+ \gets \text{VGPhrases}(M^+) \cup \text{CacheCaption}(M^+, C) \cup H^+$
\State $T^- \gets \text{VGPhrases}(M^-) \cup \text{CacheCaption}(M^-, C) \cup H^-$
\State $\vd^+_{\text{txt}} \gets \text{MeanTextEmbed}(T^+)$,\quad $\vd^-_{\text{txt}} \gets \text{MeanTextEmbed}(T^-)$
\State $\vd^+ \gets \text{Fuse}(\vd^+_{\text{img}}, \vd^+_{\text{txt}})$,\quad $\vd^- \gets \text{Fuse}(\vd^-_{\text{img}}, \vd^-_{\text{txt}})$
\State $\vq^{(t+1)} \gets \ell_2\text{-normalise}(\alpha\,\vq^{(t)} + \beta\,\vd^+ - \gamma\,\vd^-)$
\State $\mathcal{R}^{(t+1)} \gets \text{FAISSTopK}(\vq^{(t+1)}, k)$
\State \Return $\mathcal{R}^{(t+1)}$
\end{algorithmic}
\end{algorithm}

%% =============================================================================
\section{Implementation Details}
\label{sec:impl}

\subsection{Coordinate Pipeline}

User clicks are captured in CSS pixel coordinates relative to the image
display rectangle. Since the browser may scale the displayed image, the
application computes the \emph{natural pixel} coordinates of each click
by dividing by the display scale factor. These coordinates are sent to the
server as \texttt{(x, y, label)} tuples with \texttt{coord\_width} and
\texttt{coord\_height} set to the browser's \texttt{naturalWidth}/\texttt{naturalHeight}.
The server scales the clicks to the full image resolution before passing
them to SAM:
\[
  x_{\text{full}} = x_{\text{click}} \cdot W/w_{\text{coord}},
  \qquad
  y_{\text{full}} = y_{\text{click}} \cdot H/h_{\text{coord}},
\]
where $W \times H$ is the full image size and $w_{\text{coord}} \times h_{\text{coord}}$
is the natural preview size. Click coordinates are clamped to $[0, W-1] \times
[0, H-1]$ before SAM inference.

The SAM mask is produced at full resolution, then downsampled to the preview
space via nearest-neighbour interpolation. The RLE encoding uses COCO
convention (row-major, background-first, size = $[H_{\text{preview}}, W_{\text{preview}}]$).
The client decodes the RLE into a flat \texttt{Uint8ClampedArray} and renders
the mask at the correct image display rectangle by using the five-argument
form of \texttt{CanvasRenderingContext2D.drawImage}.

\subsection{SAM~3 Backbone Sharing}

SAM~3's \texttt{build\_sam3\_image\_model()} constructs the main model with a
vision-language backbone. The internal \texttt{build\_tracker()} call creates
the SAM~2-compatible interactive predictor with \texttt{with\_backbone=False},
leaving the predictor's backbone attribute as \texttt{None}. During interactive
prediction, \texttt{forward\_image} requires a non-null backbone to produce
SAM~2 feature levels.

The fix is a post-hoc attribute assignment:
\begin{verbatim}
  if self._interactive.model.backbone is None and model.backbone:
      self._interactive.model.backbone = model.backbone
\end{verbatim}
This shares the main model's backbone tensor (already on the correct device)
with the interactive predictor without duplicating memory.

\subsection{Logit Caching for Sequential Clicks}

SAM~2/3 interactive predictors support passing the logits of the previous
prediction as a mask prior, which improves mask quality when the user adds
additional click points to refine a segmentation. VisualReF caches the
last predicted logits per image path and passes them as \texttt{mask\_input}
to subsequent \texttt{predict()} calls. The cache is bounded at 50 entries
with FIFO eviction to prevent memory growth over long sessions. The logit
cache is invalidated when a negative click is added, as negative prompts
change the mask semantics.

\subsection{Asynchronous Caption Architecture}

The caption pipeline is implemented as follows. After the \texttt{/segment}
endpoint assembles its response, it calls:
\begin{verbatim}
  asyncio.create_task(_background_caption(
      image_path, label, query, region_b64, bbox, user_hint))
\end{verbatim}
The background coroutine calls Ollama in a thread pool
(\texttt{asyncio.to\_thread}), stores the result under the cache key
\texttt{f"\{path\}::\{label\}::\{query\}::\{hint\}"}, and logs success.
The \texttt{/apply\_feedback} handler receives the \texttt{\_caption\_cache}
dictionary as a parameter and checks it before invoking Ollama synchronously,
reducing the common-case captioning latency from 20--40~s to~$<$1~ms.

\subsection{Technology Stack Summary}

\begin{table}[h]
\centering
\caption{Technology stack of VisualReF.}
\label{tab:stack}
\begin{tabular}{@{}lll@{}}
\toprule
\textbf{Layer} & \textbf{Technology} & \textbf{Version / Notes} \\
\midrule
Front-end framework & Next.js + React    & 15 / 19 \\
State management    & Zustand            & 5 \\
Back-end framework  & FastAPI + Uvicorn  & 0.115 \\
Embedding model     & SigLIP large-256   & \texttt{google/siglip-large-patch16-256} \\
Segmentation        & SAM~3              & \texttt{facebook/sam3} (HF Hub) \\
VLM captioning      & Llama~3.2-Vision   & 11B, served via Ollama \\
Vector index        & FAISS FlatIP       & 1024-d, $\sim$108k vectors \\
Image corpus        & Visual Genome      & 108,249 images \\
Phrase corpus       & VG region descriptions & 4.1M region phrases \\
Device support      & CUDA / Apple MPS / CPU & PyTorch 2.x \\
Python runtime      & CPython 3.11       & \\
\bottomrule
\end{tabular}
\end{table}

%% =============================================================================
\section{Experimental Programme}
\label{sec:experiments}

The following section outlines a structured experimental programme designed to
evaluate VisualReF rigorously. Experiments are grouped by research question.

\subsection{Research Questions}

\begin{description}[leftmargin=0em,style=nextline]
  \item[\textbf{RQ1.}] Does VLM-augmented relevance feedback improve retrieval
    precision over image-only or text-only feedback?
  \item[\textbf{RQ2.}] What is the optimal fusion strategy for multimodal
    feedback signals?
  \item[\textbf{RQ3.}] How does VLM caption quality affect downstream retrieval
    performance?
  \item[\textbf{RQ4.}] Can asynchronous, progressive feedback maintain user
    satisfaction while reducing perceived latency?
  \item[\textbf{RQ5.}] How does segmentation granularity affect feedback
    quality relative to whole-image or bounding-box interaction?
  \item[\textbf{RQ6.}] What is the effect of multiple feedback rounds on query
    drift?
  \item[\textbf{RQ7.}] How does the approach scale with corpus size?
\end{description}

\subsection{Datasets}

\begin{table}[h]
\centering
\caption{Datasets used in the experimental programme.}
\label{tab:datasets}
\begin{tabular}{@{}lrll@{}}
\toprule
\textbf{Dataset}    & \textbf{Size}  & \textbf{Domain} & \textbf{Primary use} \\
\midrule
Visual Genome~\citep{Krishna2017}       & 108,249 & General       & Primary corpus, VG phrases \\
MS-COCO val2014~\citep{Lin2014}         & 40,504  & 80-cat objects & Retrieval evaluation \\
Flickr30k~\citep{Young2014}             & 31,783  & Captioned photos & Caption quality eval \\
Open Images V7 (subset)~\citep{Kuznetsova2020} & 100k--500k & Diverse & Scale experiments \\
\bottomrule
\end{tabular}
\end{table}

For retrieval evaluation on MS-COCO, relevance judgements are derived from
COCO category annotations: an image is judged relevant to a query category if
it contains at least one annotated instance of that category.

\subsection{Evaluation Metrics}

\begin{itemize}[leftmargin=1.5em]
  \item \textbf{Precision@$k$} ($k \in \{5, 10, 20\}$): fraction of top-$k$
        retrieved images judged relevant.
  \item \textbf{nDCG@$k$}: normalised discounted cumulative gain, accounting
        for ranking position.
  \item \textbf{Mean Reciprocal Rank (MRR)}: position of the first relevant
        result.
  \item \textbf{Rounds to target}: number of feedback rounds required to
        achieve P@5~$\geq$~0.8.
  \item \textbf{Wall-clock latency}: time per feedback round, broken down by
        component (SAM, SigLIP, Ollama, FAISS).
  \item \textbf{User satisfaction} (Likert, 1--5): for user study conditions.
\end{itemize}

\subsection{Baselines (RQ1)}

To isolate the contribution of each signal modality, we compare five
conditions:

\begin{enumerate}[leftmargin=1.5em]
  \item \textbf{Text-only search}: SigLIP text embedding of user query, no
        feedback.
  \item \textbf{Rocchio-Image}: SAM crops $\to$ SigLIP image embedding $\to$
        Rocchio. No text.
  \item \textbf{Rocchio-Text}: User-typed hints $\to$ SigLIP text embedding
        $\to$ Rocchio. No image signal.
  \item \textbf{Rocchio-Image+Text}: Fusion of image embeddings and user text
        embeddings. No VLM captions.
  \item \textbf{VisualReF (full)}: Image + user text + VLM captions + VG
        phrases $\to$ Rocchio. This is the proposed system.
\end{enumerate}

\subsection{Fusion Ablation (RQ2)}

We perform a grid search over fusion weights $(w_{\text{img}}, w_{\text{txt}})$
with $w_{\text{img}} \in \{0.2, 0.3, 0.4, 0.5, 0.6\}$ and $w_{\text{txt}} =
1 - w_{\text{img}}$, measuring P@5 and nDCG@10 on a held-out validation
split. Rocchio parameters $(\alpha, \beta, \gamma)$ are similarly swept:
$\alpha \in \{0.6, 0.7, 0.8, 0.9\}$, $\beta \in \{0.3, 0.5, 0.7\}$,
$\gamma \in \{0.1, 0.15, 0.2\}$.

\subsection{Caption Quality vs.\ Retrieval (RQ3)}

We compare three VLMs for crop captioning:
\begin{enumerate}[leftmargin=1.5em,noitemsep]
  \item Llama~3.2-Vision 11B (the current default).
  \item Moondream2 (1.8B, 3--5~s per crop on MPS).
  \item LLaVA-Phi3-Mini (2.7B, 5--10~s per crop).
\end{enumerate}
Caption quality is measured by BLEU-4 and ROUGE-L against Visual Genome
region descriptions as ground truth. Downstream retrieval quality is measured
by the change in P@5 between round~0 (no feedback) and round~2.

\subsection{User Study (RQ4, RQ5)}

A within-subjects user study with $n \geq 20$ participants will compare three
interaction conditions on a fixed set of 10 retrieval queries:
\begin{enumerate}[leftmargin=1.5em,noitemsep]
  \item \textbf{Whole-image}: users mark full images as relevant/irrelevant.
  \item \textbf{Bounding-box}: users draw bounding boxes around regions of
        interest.
  \item \textbf{SAM point-click}: the VisualReF interaction mode.
\end{enumerate}
Within condition~(c), we additionally vary synchronous (blocking on VLM)
vs.\ asynchronous (background caption, immediate preliminary results) to
address RQ4. Measures include task completion time, number of feedback
rounds, P@5 at round 3, and a post-task questionnaire on satisfaction,
perceived responsiveness, and ease of use.

\subsection{Query Drift Analysis (RQ6)}

For 50 queries across 10 feedback rounds each, we track the cosine similarity
between $\vq^{(t)}$ and $\vq^{(0)}$ and correlate it with P@5 to characterise
drift dynamics. We compare three conditions: no drift mitigation, fuse-initial
($\vq_{\text{fused}}^{(t)}$ anchor), and reset to $\vq^{(0)}$ every~3~rounds.

\subsection{Scale Experiments (RQ7)}

FAISS indexes are built at five corpus sizes: 10k, 40k (COCO), 108k (VG),
250k, and 500k images (sampled from Open Images). At each scale we report:
P@5 at round~2, mean query latency, and index memory footprint. We additionally
compare \texttt{IndexFlatIP} (exact) against \texttt{IndexHNSWFlat} (16 links)
for the 108k and 500k settings.

%% =============================================================================
\section{Discussion}
\label{sec:discussion}

\subsection{Latency and the VLM Bottleneck}

The primary engineering challenge in VisualReF is VLM inference latency. On
Apple MPS, Llama~3.2-Vision requires 20--40~seconds per crop. The asynchronous
captioning architecture mitigates this by hiding the latency behind the user's
annotation time: the caption is computed in the background while the user
inspects results, marks regions, and types hints. In practice, when the user
takes more than 30~seconds between the first SAM click and pressing ``Apply
Feedback'', the cache is warm and \texttt{/apply\_feedback} returns in under
2~seconds.

This architectural pattern—\emph{speculative VLM inference}—is a practical
contribution applicable to other systems that combine slow neural inference
with interactive user feedback.

\subsection{Embedding Space Alignment}

The effectiveness of the fusion scheme depends on SigLIP's ability to align
image and text embeddings of semantically equivalent content. While SigLIP was
trained for global image-text alignment, the feedback signals operate on
\emph{cropped regions}, which may differ substantially from the global image
that a SigLIP caption would naturally describe. It is an empirical question
whether SigLIP region embeddings are sufficiently aligned with region-level
text descriptions. The experimental programme in \cref{sec:experiments}
is designed to test this assumption.

\subsection{Visual Genome Region Phrases as a Proxy}

The use of VG region descriptions as a fallback for VLM captions is
pragmatically motivated, but the proxy is imperfect: VG descriptions were
written by crowdworkers following standardised instructions, whereas the VLM
captions are conditioned on the user's specific query. When the query is
\emph{``dog playing near water at sunset''} and the VG phrase for the selected
region is simply \emph{``a dog''}, the phrase provides much weaker signal.
The experimental comparison between VG-only and VLM-augmented conditions (RQ1)
will quantify this gap.

\subsection{Rocchio Limitations}

The Rocchio update is a linear operation in a fixed embedding space. It assumes
that the relevant and irrelevant feedback signals can be combined additively,
which may not hold when the relevant and irrelevant subspaces are entangled in
the SigLIP embedding space. Non-linear query update methods—such as neural
query reformulation~\citep{Jiang2015} or learned feedback
encoders~\citep{Yu2021}—may achieve better performance at the cost of
additional model complexity. The present work treats Rocchio as the baseline
and does not explore these alternatives.

\subsection{Limitations of the Prototype}

\begin{itemize}[leftmargin=1.5em]
  \item \textbf{Single-process back-end:} The FastAPI server runs in a single
        Python process. SAM inference serialises all user clicks through a
        mutex (\texttt{asyncio.Lock}), preventing concurrent segmentation.
  \item \textbf{No persistent sessions:} Query state is held in server memory
        and lost on restart. A production system would require session storage.
  \item \textbf{FAISS IndexFlatIP:} Exact brute-force search does not scale
        beyond~$\sim$500k images at acceptable latency.
  \item \textbf{Fixed Rocchio parameters:} The parameters $(\alpha, \beta,
        \gamma)$ are constant across queries and users. Adaptive parameters
        may better serve diverse query types.
  \item \textbf{No evaluation against human relevance judgements:}
        The current prototype has not yet been evaluated in a controlled user
        study. All claims about retrieval effectiveness are based on proxy
        metrics (category overlap) rather than direct human assessment.
\end{itemize}

%% =============================================================================
\section{Conclusion}
\label{sec:conclusion}

This thesis has presented VisualReF, a prototype interactive image retrieval
system that combines SAM~3 interactive segmentation, SigLIP vision-language
embeddings, and Llama~3.2-Vision VLM captioning into a multimodal Rocchio
relevance feedback loop. The key technical contributions are: a formalisation
of the multimodal feedback update, an asynchronous captioning architecture
that hides VLM latency, and a hybrid signal strategy that leverages Visual
Genome region descriptions as instant textual feedback when VLM inference is
unavailable.

The system demonstrates that the semantic gap in visual relevance feedback
can be bridged by a three-way signal combination: \emph{what the user sees}
(SAM crop image embedding), \emph{what the VLM observes} (query-conditioned
caption embedding), and \emph{what the user articulates} (free-text hint
embedding). The asynchronous architecture makes this combination feasible at
interactive latency despite the high cost of large VLM inference.

The experimental programme outlined in \cref{sec:experiments} will test the
central hypothesis that this multimodal combination improves retrieval
precision over image-only or text-only feedback, and will quantify the
trade-offs between caption quality, latency, and retrieval effectiveness.

\textbf{Future work} includes: adapting the Rocchio parameters online using a
learned feedback model, replacing the synchronous Ollama call with a
distilled region-description model fine-tuned for SigLIP alignment, scaling
the FAISS index to the full Open Images dataset, and conducting a formal
user study to evaluate the interactive experience.

%% =============================================================================
\section*{Acknowledgements}

The author thanks the supervisors at Maastricht University for their guidance,
and the original VisualReF demo authors—Bulat Khaertdinov, Mirela Popa, and
Nava Tintarev—for the foundational prototype on which this work builds.
SAM~3 is provided under Meta's open model licence; Visual Genome is released
under Creative Commons Attribution 4.0; all other model weights are used under
their respective HuggingFace Hub licences.

%% =============================================================================
%% Bibliography
%% =============================================================================
\bibliographystyle{plainnat}
\begin{thebibliography}{99}

\bibitem[Babenko et al.(2014)]{Babenko2014}
Babenko, A., Slesarev, A., Chigorin, A., and Lempitsky, V. (2014).
Neural codes for image retrieval.
In \emph{Proceedings of ECCV 2014}, pages 584--599. Springer.

\bibitem[Boykov and Jolly(2001)]{Boykov2001}
Boykov, Y. and Jolly, M.-P. (2001).
Interactive graph cuts for optimal boundary and region segmentation of objects in N-D images.
In \emph{Proceedings of ICCV 2001}, pages 105--112. IEEE.

\bibitem[Custodio and Sebe(2023)]{Custodio2023}
Custodio, E. and Sebe, N. (2023).
Leveraging language for accelerated learning of tool use.
\emph{arXiv preprint arXiv:2309.13977}.

\bibitem[Flickner et al.(1995)]{Flickner1995}
Flickner, M., Sawhney, H., Niblack, W., Ashley, J., Huang, Q., Dom, B., Gorkani, M.,
Hafner, J., Lee, D., Petkovic, D., Steele, D., and Yanker, P. (1995).
Query by image and video content: The QBIC system.
\emph{IEEE Computer}, 28(9):23--32.

\bibitem[Gosselin and Cord(2004)]{Gosselin2004}
Gosselin, P.-H. and Cord, M. (2004).
Active learning methods for interactive image retrieval.
\emph{IEEE Transactions on Image Processing}, 17(7):1200--1211.

\bibitem[Jia et al.(2021)]{Jia2021}
Jia, C., Yang, Y., Xia, Y., Chen, Y.-T., Parekh, Z., Pham, H., Le, Q., Sung, Y.,
Li, Z., and Duerig, T. (2021).
Scaling up visual and vision-language representation learning with noisy text supervision.
In \emph{Proceedings of ICML 2021}, pages 4904--4916. PMLR.

\bibitem[Jiang et al.(2015)]{Jiang2015}
Jiang, J., He, D., and Lin, J. (2015).
Automatic online news topic ranking using media focus and user attention based on
implicit link analysis.
In \emph{Proceedings of CIKM 2015}.

\bibitem[Johnson et al.(2019)]{Johnson2019}
Johnson, J., Douze, M., and Jégou, H. (2019).
Billion-scale similarity search with GPUs.
\emph{IEEE Transactions on Big Data}, 7(3):535--547.

\bibitem[Khaertdinov et al.(2025)]{Khaertdinov2025}
Khaertdinov, B., Popa, M., and Tintarev, N. (2025).
{VisualReF}: Interactive image search prototype with visual relevance feedback.
In \emph{Proceedings of ACM RecSys~'25}. ACM.
\texttt{doi:10.1145/3705328.3759341}

\bibitem[Kirillov et al.(2023)]{Kirillov2023}
Kirillov, A., Mintun, E., Ravi, N., Mao, H., Rolland, C., Gustafson, L., Xiao, T.,
Whitehead, S., Berg, A. C., Lo, W.-Y., Dollar, P., and Girshick, R. (2023).
Segment anything.
In \emph{Proceedings of ICCV 2023}. IEEE.

\bibitem[Krishna et al.(2017)]{Krishna2017}
Krishna, R., Zhu, Y., Groth, O., Johnson, J., Hata, K., Kravitz, J., Chen, S.,
Kalantidis, Y., Li, L.-J., Shamma, D. A., Bernstein, M. S., and Fei-Fei, L. (2017).
Visual genome: Connecting language and vision using crowdsourced dense image annotations.
\emph{International Journal of Computer Vision}, 123(1):32--73.

\bibitem[Kuznetsova et al.(2020)]{Kuznetsova2020}
Kuznetsova, A., Rom, H., Alldrin, N., Uijlings, J., Krasin, I., Pont-Tuset, J.,
Kamali, S., Popov, S., Malloci, M., Kolesnikov, A., Duerig, T., and Ferrari, V. (2020).
The Open Images dataset V4.
\emph{International Journal of Computer Vision}, 128(7):1956--1981.

\bibitem[Lin et al.(2014)]{Lin2014}
Lin, T.-Y., Maire, M., Belongie, S., Hays, J., Perona, P., Ramanan, D., Dollar, P.,
and Zitnick, C. L. (2014).
Microsoft COCO: Common objects in context.
In \emph{Proceedings of ECCV 2014}, pages 740--755. Springer.

\bibitem[Liu et al.(2023)]{Liu2023}
Liu, H., Li, C., Wu, Q., and Lee, Y. J. (2023).
Visual instruction tuning.
In \emph{Proceedings of NeurIPS 2023}. Curran Associates.

\bibitem[Luo et al.(2023)]{Luo2023}
Luo, H., Zhang, X., Deng, C., Wang, X., and Bao, H. (2023).
Cheap and quick: Efficient vision-language instruction tuning for large language models.
\emph{arXiv preprint arXiv:2305.15023}.

\bibitem[Meta(2024)]{Meta2024}
Meta AI. (2024).
Llama 3.2: Revolutionising edge AI and vision.
\url{https://ai.meta.com/blog/llama-3-2-connect-2024-vision-edge-mobile-devices/}

\bibitem[Meta(2025)]{Meta2025}
Meta AI. (2025).
SAM 3: Segment Anything Model version 3.
\url{https://ai.meta.com/research/publications/sam3/}

\bibitem[Radford et al.(2021)]{Radford2021}
Radford, A., Kim, J. W., Hallacy, C., Ramesh, A., Goh, G., Agarwal, S., Sastry, G.,
Askell, A., Mishkin, P., Clark, J., Krueger, G., and Sutskever, I. (2021).
Learning transferable visual models from natural language supervision.
In \emph{Proceedings of ICML 2021}, pages 8748--8763. PMLR.

\bibitem[Ravi et al.(2024)]{Ravi2024}
Ravi, N., Gabeur, V., Hu, Y.-T., Hu, R., Ryali, C., Ma, T., Khedr, H., R{\"a}dle, R.,
Rolland, C., Gustafson, L., Mintun, E., Pan, J., Alwala, K. V., Carion, N.,
Wu, C.-Y., Girshick, R., Dollar, P., and Feichtenhofer, C. (2024).
SAM 2: Segment anything in images and videos.
\emph{arXiv preprint arXiv:2408.00714}.

\bibitem[Rocchio(1971)]{Rocchio1971}
Rocchio, J. J. (1971).
Relevance feedback in information retrieval.
In Salton, G. (ed.), \emph{The SMART Retrieval System: Experiments in Automatic
Document Processing}, pages 313--323. Prentice-Hall.

\bibitem[Rother et al.(2004)]{Rother2004}
Rother, C., Kolmogorov, V., and Blake, A. (2004).
GrabCut: Interactive foreground extraction using iterated graph cuts.
\emph{ACM Transactions on Graphics (SIGGRAPH)}, 23(3):309--314.

\bibitem[Rui et al.(1998)]{Rui1998}
Rui, Y., Huang, T. S., Ortega, M., and Mehrotra, S. (1998).
Relevance feedback: A power tool for interactive content-based image retrieval.
\emph{IEEE Transactions on Circuits and Systems for Video Technology}, 8(5):644--655.

\bibitem[Smeulders et al.(2000)]{Smeulders2000}
Smeulders, A. W. M., Worring, M., Santini, S., Gupta, A., and Jain, R. (2000).
Content-based image retrieval at the end of the early years.
\emph{IEEE Transactions on Pattern Analysis and Machine Intelligence}, 22(12):1349--1380.

\bibitem[Smith and Chang(1996)]{Smith1996}
Smith, J. R. and Chang, S.-F. (1996).
VisualSEEk: A fully automated content-based image query system.
In \emph{Proceedings of ACM Multimedia 1996}.

\bibitem[Sofiiuk et al.(2022)]{Sofiiuk2022}
Sofiiuk, K., Petrov, I., and Konushin, A. (2022).
Reviving iterative training with mask guidance for interactive segmentation.
In \emph{Proceedings of ICIP 2022}. IEEE.

\bibitem[Tolias et al.(2016)]{Tolias2016}
Tolias, G., Sicre, R., and Jégou, H. (2016).
Particular object retrieval with integral max-pooling of CNN activations.
In \emph{Proceedings of ICLR 2016}.

\bibitem[Tong and Chang(2001)]{Tong2001}
Tong, S. and Chang, E. (2001).
Support vector machine active learning for image retrieval.
In \emph{Proceedings of ACM Multimedia 2001}, pages 107--118.

\bibitem[Xu et al.(2016)]{Xu2016}
Xu, N., Price, B., Cohen, S., Yang, J., and Huang, T. (2016).
Deep interactive object selection.
In \emph{Proceedings of CVPR 2016}. IEEE.

\bibitem[Young et al.(2014)]{Young2014}
Young, P., Lai, A., Hodosh, M., and Hockenmaier, J. (2014).
From image descriptions to visual denotations: New similarity metrics for semantic
inference over event descriptions.
\emph{Transactions of the Association for Computational Linguistics}, 2:67--78.

\bibitem[Yu et al.(2021)]{Yu2021}
Yu, T., Meng, J., Cho, Y., Cao, J., and Yuan, J. (2021).
Multi-query interactive image retrieval with relevance feedback.
In \emph{Proceedings of IJCAI 2021}.

\bibitem[Zhai et al.(2023)]{Zhai2023}
Zhai, X., Kolesnikov, A., Houlsby, N., and Beyer, L. (2023).
Sigmoid loss for language image pre-training.
In \emph{Proceedings of ICCV 2023}. IEEE.

\bibitem[Zhou and Huang(2003)]{Zhou2003}
Zhou, X. S. and Huang, T. S. (2003).
Relevance feedback in image retrieval: A comprehensive review.
\emph{Multimedia Systems}, 8(6):536--544.

\end{thebibliography}

\end{document}
